% ==========================================
% FILE: Full_Cover.tex (ĐỘC LẬP HOÀN TOÀN)
% ==========================================
\documentclass{article}
\usepackage{tikz}
\usetikzlibrary{calc, positioning}
\usepackage{xcolor}

% --- BƯỚC 1: KHAI BÁO BIẾN SỐ (CSS VARIABLES) ---
% Hãy thay đổi các con số này, toàn bộ thiết kế sẽ tự động co giãn!
\pgfmathsetmacro{\BookW}{13}     % Bề ngang 1 trang (cm) - Chuẩn B6
\pgfmathsetmacro{\BookH}{19}     % Bề dọc 1 trang (cm)
\pgfmathsetmacro{\Spine}{1.8}    % Độ dày GÁY SÁCH (cm) - Tính theo số trang
\pgfmathsetmacro{\Flap}{8}       % Độ rộng TAY GẤP (cm)
\pgfmathsetmacro{\Bleed}{0.3}    % Lề xén/Bleed (cm) - Chuẩn xưởng in

% Tính toán tự động tổng kích thước (Cộng dồn)
\pgfmathsetmacro{\TotalW}{\Bleed + \Flap + \BookW + \Spine + \BookW + \Flap + \Bleed}
\pgfmathsetmacro{\TotalH}{\Bleed + \BookH + \Bleed}

% Ép khổ giấy PDF đúng bằng tổng kích thước vừa tính
\usepackage[
    paperwidth=\TotalW cm, 
    paperheight=\TotalH cm, 
    margin=0pt
]{geometry}

% Màu sắc
\definecolor{coverbg}{cmyk}{0, 0.05, 0.1, 0} % Màu nền kem nhẹ
\definecolor{foldline}{cmyk}{0, 1, 1, 0}     % Màu đỏ (Chỉ để nhìn vạch gấp, lúc in sẽ ẩn đi)

\begin{document}
\thispagestyle{empty}
\pagecolor{coverbg}

\begin{tikzpicture}[remember picture, overlay]

    % --- BƯỚC 2: KHAI BÁO CÁC ĐIỂM CHỐT (CSS GRID) ---
    % Đặt gốc (0,0) tại góc dưới cùng bên trái của trang giấy
    \coordinate (Origin) at (current page.south west);

    % Tính toán tọa độ X của các đường gấp
    \pgfmathsetmacro{\XFlapLeft}{\Bleed + \Flap}
    \pgfmathsetmacro{\XSpineLeft}{\XFlapLeft + \BookW}
    \pgfmathsetmacro{\XSpineRight}{\XSpineLeft + \Spine}
    \pgfmathsetmacro{\XFlapRight}{\XSpineRight + \BookW}

    % Vẽ các đường đứt nét mô phỏng NẾP GẤP (Bản in thật hãy comment/ẩn phần này đi)
    \draw[foldline, dashed] ($(Origin) + (\XFlapLeft, 0)$) -- ($(Origin) + (\XFlapLeft, \TotalH)$);
    \draw[foldline, dashed] ($(Origin) + (\XSpineLeft, 0)$) -- ($(Origin) + (\XSpineLeft, \TotalH)$);
    \draw[foldline, dashed] ($(Origin) + (\XSpineRight, 0)$) -- ($(Origin) + (\XSpineRight, \TotalH)$);
    \draw[foldline, dashed] ($(Origin) + (\XFlapRight, 0)$) -- ($(Origin) + (\XFlapRight, \TotalH)$);

    % --- BƯỚC 3: ĐẶT NỘI DUNG VÀO GIỮA CÁC KHUNG (FLEX CENTER) ---

    % 1. Tay gấp Trái (Tóm tắt truyện)
    \node[text width=\Flap cm - 2cm, align=justify, anchor=center]
    at ($(Origin) + (\Bleed + \Flap/2, \TotalH/2)$) {
        \sffamily\small
        \textbf{TÓM TẮT}\\
        Mười năm trước, cánh cổng ma thuật mở ra... (Chữ tự động Wrap theo khung tay gấp).
    };

    % 2. Bìa Sau (Back Cover - Ảnh & Barcode)
    \node[anchor=center]
    at ($(Origin) + (\XFlapLeft + \BookW/2, \TotalH/2)$) {
        \Huge \textbf{BÌA SAU}
    };

    % 3. Gáy Sách (Spine - Xoay ngang 90 độ)
    \node[anchor=center, rotate=-90]
    at ($(Origin) + (\XSpineLeft + \Spine/2, \TotalH/2)$) {
        \Large \textbf{TÊN TRUYỆN — TẬP 1}
    };

    % 4. Bìa Trước (Front Cover - Tên truyện & Tác giả)
    \node[anchor=center]
    at ($(Origin) + (\XSpineRight + \BookW/2, \TotalH/2)$) {
        \begin{tabular}{c}
            \Huge \textbf{TÊN TRUYỆN} \\
            \vspace{1cm}              \\
            \Large \textit{Tác giả: Akihiko Sora}
        \end{tabular}
    };

    % 5. Tay gấp Phải (Giới thiệu Tác giả)
    \node[text width=\Flap cm - 2cm, align=center, anchor=center]
    at ($(Origin) + (\XFlapRight + \Flap/2, \TotalH/2)$) {
        \sffamily\small
        \textbf{TÁC GIẢ}\\
        Một người thích FOSS và đam mê toán học.
    };

\end{tikzpicture}
\end{document}